\documentclass[11pt]{article}
\usepackage[T1]{fontenc}
\usepackage[utf8]{inputenc}
\usepackage{lmodern,amsmath,amssymb,amsthm,mathtools}
\usepackage[margin=28mm]{geometry}
\usepackage{microtype,xcolor}
\usepackage[authoryear,round]{natbib}
\usepackage{xurl}
\usepackage[colorlinks=true,allcolors=blue!50!black]{hyperref}
\newtheorem{theorem}{Theorem}[section]
\newtheorem{proposition}[theorem]{Proposition}
\newtheorem{lemma}[theorem]{Lemma}
\newtheorem{corollary}[theorem]{Corollary}
\theoremstyle{definition}\newtheorem{definition}[theorem]{Definition}
\theoremstyle{remark}\newtheorem{remark}[theorem]{Remark}
\newcommand{\OO}{\mathcal O}
\newcommand{\PP}{\mathbf P}
\newcommand{\FF}{\mathbf F}
\newcommand{\Pic}{\operatorname{Pic}}
\newcommand{\NS}{\operatorname{NS}}
\newcommand{\End}{\operatorname{End}}
\newcommand{\Hom}{\operatorname{Hom}}
\newcommand{\Nm}{\operatorname{Nm}}
\newcommand{\Tr}{\operatorname{Tr}}
\newcommand{\ann}{\operatorname{ann}}
\newcommand{\rk}{\operatorname{rank}}
\newcommand{\diag}{\operatorname{diag}}
\newcommand{\workbenchurl}{https://github.com/KokunoYumeto/schemes-and-moduli-workbench}
\newcommand{\recordurl}{https://github.com/KokunoYumeto/mathematics-workbenches}
\title{Grothendieck's generic Frobenius-rank question for abelian surfaces in characteristic two}
\author{}
\date{23 November 2026}
\begin{document}
\maketitle
\begin{abstract}
Grothendieck's crystalline-cohomology notes relate generic stable Frobenius ranks of hypersurface sections to unit roots and ask for eventual positivity. For an abelian surface $A$ in characteristic two and an ample $M$ of odd Euler characteristic $D$, we determine the generic Frobenius-stable dimension of the relative coherent cohomology of curves in $|M^{3d}|$, for every positive integer $d$. It is $9Dd^2-1$ in odd degree and $9Dd^2-3+f_A$ in even degree, where $f_A$ is the $2$-rank of $A$. Thus the generic relative Frobenius is bijective in every odd degree, whereas every even degree retains a defect when $A$ is nonordinary. For an ordinary surface and an arbitrary ample $M$, a generic curve in $|M^n|$ is ordinary for every $n\ge3$. The proof combines characteristic-two theta jets, a nine-dimensional Cartier tensor, the Pfaffian cubic of Gruson and Sam, Oda's simplicity theorem, finite isotropic descent, nodal norm images, and an integral Prym bound. All-degree assertions concern complete linear systems; the precise additional condition for ambient hypersurfaces is stated. This text and its checks were produced by GPT-6 Astra Pro; the extent of the independent checks is specified at the end, and no human has checked the mathematics.
\end{abstract}

\section{The rank formulas}

The stable Frobenius rank records the part of coherent cohomology that survives repeated Frobenius. For hypersurface sections it controls the unit-root contribution in Grothendieck's crystalline-cohomology programme. We determine it in complete linear systems on two large classes of abelian surfaces, including every principally polarized surface in characteristic two. The resulting parity formula shows that an asymptotic positivity statement can coexist with infinitely many degrees of nonbijective relative Frobenius.

Work over an algebraically closed field $k$ of characteristic two, and let $A$ be an abelian surface over $k$. Results over a perfect field mean the corresponding geometric generic results after algebraic closure. For an abelian variety $B$, its $2$-rank $f_B$ is determined by $|B[2](k)|=2^{f_B}$. For a smooth connected curve $C$, put $f_C=f_{\operatorname{Jac}(C)}$. Equivalently, $f_C$ is the dimension of the largest Frobenius-bijective subspace of $H^1(C,\OO_C)$.

For a smooth ample curve $C\subset A$ define
\[
 Q_C=\operatorname{coker}\bigl(H^1(A,\OO_A)\longrightarrow H^1(C,\OO_C)\bigr),
 \qquad r(C)=\dim_k(Q_C)_{\mathrm{st}}.
\]
Here $V_{\mathrm{st}}=\bigcap_{j\ge0}\operatorname{im}F^j$ for a finite-dimensional space with semilinear Frobenius. The map defining $Q_C$ is injective. If $C\in|M^n|$ and $D=\chi(M)$, then
\begin{equation}\label{eq:dimensions}
 g(C)=Dn^2+1,\qquad \dim Q_C=Dn^2-1,\qquad r(C)=f_C-f_A.
\end{equation}
The preliminary lemmas below prove these assertions and the openness needed to speak of a generic rank.

\begin{theorem}[Odd polarization degree]\label{thm:odd}
Let $M$ be ample on $A$ and suppose that $D=\chi(M)$ is odd. For each $d\ge1$, a nonempty open of smooth curves in $|M^{3d}|$ has relative stable rank
\[
 r_d=\begin{cases}
 9Dd^2-1,&d\text{ odd},\\
 9Dd^2-3+f_A,&d\text{ even}.
 \end{cases}
\]
The corresponding full curve defect $g(C)-f_C$ is $2-f_A$ in odd degree and $2(2-f_A)$ in even degree. These are the generic ranks of the complete linear systems.
\end{theorem}

For a principal polarization $\Theta$, take $M=\OO_A(\Theta)$ and $D=1$. The result applies to every principally polarized surface, including products and supersingular surfaces, without any restriction on the field of definition of its moduli point.

\begin{theorem}[Ordinary surfaces]\label{thm:ordinary}
Suppose $A$ is ordinary. For every ample line bundle $M$ and every $n\ge3$, a generic smooth curve in $|M^n|$ is ordinary. In particular, with $D=\chi(M)$,
\[
 f_C=g(C)=Dn^2+1,\qquad r(C)=Dn^2-1.
\]
No condition on the parity of $D$ or $n$ is required.
\end{theorem}

\begin{corollary}\label{cor:parity}
On a nonordinary principally polarized surface, the generic relative Frobenius in $|3d\Theta|$ is bijective exactly when $d$ is odd. Its stable rank is positive in every degree, but there is no threshold beyond which it is bijective in every degree.
\end{corollary}

For example, if $D=1$ and $f_A=0$, then the first cubic system has $g(C)=10$, $f_C=8$ and $r(C)=8=\dim Q_C$. In the next system $|6\Theta|$, the genus is $37$, the full $2$-rank is $33$, and the relative space has dimension $35$ but stable rank $33$. The even-degree loss is already visible in these two smallest degrees.

In \S8.5.8 of his crystalline-cohomology notes, Grothendieck considers the maximal generic semisimple rank for hypersurface sections and asks for its eventual positivity. The continuation carries the marginal comment ``Or cela est \'etabli au n\textsuperscript{o} 7'' \citep[cote~4, archive sheets~7--8, PDF images~8--9]{Grothendieck}. The formulas above answer a more precise rank question for the indicated abelian surfaces. They distinguish positivity, which holds here in every positive degree, from bijectivity, which exhibits an enduring parity obstruction. They do not give a rank formula for an arbitrary ambient variety.

The inputs concerning theta groups and polarizations are classical; exact statements are cited below. The simplicity input is \citet[Theorem~1.2(i), p.~45]{Oda}, whose Section~1 treats abelian varieties of arbitrary dimension and allows inseparable isogenies. The induced Jacobian polarization is computed by \citet[\S3, equations~(3.2)--(3.3), pp.~346--347]{Debarre}. The four-parameter Pfaffian cubic is already explicit in \citet[equation~(3.9), author version p.~14; \S3.4, pp.~17--18]{GrusonSam}, together with its forty-point parameter base locus. Their geometric treatment assumes characteristic different from two and three. We reprove the polynomial identity and the needed exceptional-tensor decomposition in characteristic two, then apply them to the Cartier tensor. \citet[Theorem~5.1, p.~340; Remark~5.8, p.~342]{RainsSam} extend the relation between genus-two geometry and trivectors to arbitrary fields and describe a characteristic-two Pfaffian cubic in the branch-locus construction. That construction provides neighboring geometric context; no identification with our Cartier tensor is assumed. The present all-degree arguments use the relative Frobenius determinant, nodal norm images and the integral Prym bound. A bounded literature search did not locate these exact rank formulas; this observation is not a priority claim.

The main difficulty is to turn a first Cartier calculation into a stable rank and then preserve a sharp defect in every degree. For a cubic principal section, the Heisenberg action makes the Cartier tensor symmetric. In the nonordinary case its Pfaffian cubic could a priori vanish; the exceptional tensors would produce a nontrivial projector on the globally generated bundle $F_*\OO(3\Theta)$, which Oda's simplicity theorem forbids. An explicit determinant then proves bijectivity on the relative quotient itself. Finite isotropic descent supplies principal norm sources: odd polarization degree makes their kernels prime to two, while ordinarity supplies an alternative multiplicative isotropic subgroup whose Cartier dual is \'etale. Nodal norm images preserve the normalization defect, and unions give the even degrees. An integral polarization identity forces the matching even-degree Prym bound. The statements concern complete systems; the ambient-hypersurface transfer and its restriction-surjectivity condition are given in Section~7. Other nonordinary even-degree polarization classes and general ambient varieties remain outside these theorems.

\section{Cohomology, stable rank, and nodal curves}

We use the usual polarization homomorphism
\[
 \lambda_N:A\longrightarrow A^\vee=\Pic^0(A),\qquad
 a\longmapsto t_a^*N\otimes N^{-1}.
\]
For an ample $N$, $\lambda_N$ is an isogeny, $\deg\lambda_N=\chi(N)^2$, and $\chi(N)=N^2/2$ on a surface. A numerical equivalence of line bundles is denoted by $\equiv$. On an abelian variety it is equivalence modulo $\Pic^0$, and $\NS(A)$ is torsion-free. Indeed $[N]\mapsto\lambda_N$ embeds it into the torsion-free group $\Hom(A,A^\vee)$; see \citet[\S7 and Proposition~11.2--Corollary~11.3, pp.~159--160]{EGM}.

\begin{lemma}\label{lem:cohomology}
If $N$ is ample on an abelian surface, then $H^i(A,N)=0$ for $i>0$ and $H^1(A,N^{-1})=0$. For every effective divisor $C\in|N|$, the curve is connected, its arithmetic genus is $\chi(N)+1$, and there is a Frobenius-compatible exact sequence
\begin{equation}\label{eq:relative}
 0\longrightarrow H^1(A,\OO_A)\longrightarrow H^1(C,\OO_C)
 \longrightarrow H^2(A,N^{-1})\longrightarrow H^2(A,\OO_A)
 \longrightarrow0.
\end{equation}
On $H^2(A,N^{-1})$, the semilinear operator in this sequence is the composite
\[
 H^2(A,N^{-1})\xrightarrow{F}H^2(A,N^{-2})\xrightarrow{\,s\,}H^2(A,N^{-1}).
\]
The last map of the sequence is dual to the inclusion $k s\subset H^0(A,N)$, where $s$ defines $C$. In particular its kernel has dimension $\chi(N)-1$.
\end{lemma}
\begin{proof}
An ample bundle has a symmetric translate. To see this, write $[-1]^*N\otimes N^{-1}=P\in\Pic^0(A)$, choose $Q$ with $Q^2=P$, and observe that $N\otimes Q$ is symmetric. Surjectivity of $\lambda_N$ realizes this twist as a translate. For a symmetric representative and an arbitrarily large odd integer $u$, one has $[u]^*N=N^{u^2}$. The trace of the finite \'etale map $[u]$, whose degree is the invertible scalar $u^4$, splits the map on structure sheaves. Projection formula therefore injects $H^i(A,N)$ into $H^i(A,N^{u^2})$. Serre vanishing proves the first assertion; duality and $\omega_A\simeq\OO_A$ prove the second. Translation gives the assertion for the original bundle and all its algebraically trivial twists.

Apply cohomology to $0\to N^{-1}\xrightarrow{s}\OO_A\to\OO_C\to0$. Since $H^0(N^{-1})=H^1(N^{-1})=0$, one obtains $H^0(\OO_C)=k$ and the displayed sequence. Its terminal surjectivity follows either from $H^2(C,\OO_C)=0$ or from duality. The dual map sends $1$ to $s$. Riemann--Roch gives $h^0(N)=\chi(N)$ and adjunction gives $p_a(C)=1+N^2/2$. Frobenius compatibility follows from the restriction-sequence map whose left-hand operation sends a local section $u$ of $N^{-1}$ to $s u^2$: multiplication by the defining section then gives $s^2u^2=(su)^2$ in $\OO_A$.
\end{proof}

\begin{lemma}\label{lem:stable}
Stable parts are exact on finite-dimensional semilinear Frobenius spaces over $k$. In a proper flat family of connected divisors from a fixed complete linear system on $A$, the locus where the stable rank of $H^1(\OO_C)$ is at least a specified integer is open. The same assertion holds for the relative quotient on smooth fibres.
\end{lemma}
\begin{proof}
For a space of dimension $b$, the images of iterated Frobenius stabilize after at most $b$ steps. Since $k$ is perfect, their dimensions behave exactly as for linear maps. The stabilized image is bijective under Frobenius and its complementary generalized kernel is nilpotent; the direct-sum decomposition follows by the usual Fitting argument applied to a sufficiently high iterate. A map respects both parts. Given a surjection, lift a vector in the stable part and repeatedly apply Frobenius after inverting it on the stable target; the nilpotent components disappear, proving surjectivity on stable parts. Injectivity and exactness in the middle follow from the same decomposition.

The preceding lemma makes $h^0(\OO_C)=1$ and $h^1(\OO_C)=\chi(N)+1$ constant throughout the divisor family. Cohomology and base change give a vector bundle carrying a linearized Frobenius map. In a local basis with matrix $B$, the $b$th iterate has matrix $B B^{[2]}\cdots B^{[2^{b-1}]}$, with the order depending only on the chosen linearization convention. Its rank is the stable rank on a geometric fibre. Nonvanishing of a minor is open. The relative spaces are the constant-rank kernels in \eqref{eq:relative}, so the same argument applies there. Exactness of stable parts gives \eqref{eq:dimensions}, using the standard cohomological description of the $2$-rank.
\end{proof}

\begin{lemma}[Defect of a nodal curve]\label{lem:nodal}
Let $C$ be a connected reduced projective nodal curve with smooth normalization components $C_1,\ldots,C_v$. If its dual graph has first Betti number $b_1$, then
\[
 h^1(C,\OO_C)=\sum_i g(C_i)+b_1,\qquad
 \dim H^1(C,\OO_C)_{\mathrm{st}}=\sum_i f_{C_i}+b_1.
\]
Thus its defect is the sum $\sum_i(g(C_i)-f_{C_i})$.
\end{lemma}
\begin{proof}
The normalization sequence is
\[
 0\to\OO_C\to\bigoplus_i\nu_*\OO_{C_i}\to\bigoplus_{x\text{ node}}k_x\to0,
\]
where the last map is the difference of the two branch values. The induced map on constants is the graph coboundary. Its cokernel is $k^{b_1}$ and gives an exact sequence
$0\to k^{b_1}\to H^1(C,\OO_C)\to\bigoplus_iH^1(C_i,\OO_{C_i})\to0$.
Frobenius is coordinatewise squaring on the graph term and hence is bijective. Apply Lemma~\ref{lem:stable}.
\end{proof}

We will construct nodal members with a given defect and use openness to produce smooth members of no larger defect. Very ampleness supplies a nonempty smooth-member open; irreducibility of the complete linear system makes these two opens meet.

\section{Theta divisors and jets in characteristic two}

\begin{lemma}\label{lem:dichotomy}
A principally polarized abelian surface has either a product theta divisor on a product of elliptic curves or a smooth genus-two theta divisor. In the second case the polarized surface is the Jacobian of that curve.
\end{lemma}
\begin{proof}
Write an effective theta divisor as $\Theta=\sum m_iB_i$. Abelian varieties contain no rational curves: after removing any inseparable part from a putative map from $\PP^1$, a nonzero differential would pull invariant differentials back to nonzero regular differentials on $\PP^1$. For every integral $B_i$, adjunction therefore gives $B_i^2\ge0$, with this integer even. Since $\Theta^2=2$ and every $\Theta B_i$ is positive, reducibility forces precisely two reduced components. A divisor $mB$ with $m>1$ is impossible because $m^2B^2=2$. Their equations $B_i^2+B_1B_2=1$ imply $B_i^2=0$ and $B_1B_2=1$. They are smooth elliptic translates, and addition of the corresponding elliptic subgroups is an isogeny of degree one. This identifies the polarization with the product polarization.

If $\Theta$ is integral, its arithmetic genus is two. Its normalization cannot have genus one, since an elliptic curve maps to an elliptic translate, of self-intersection zero, on an abelian variety. Thus $\Theta$ is smooth of genus two. Its inclusion induces a surjection $u:J(\Theta)\to A$ and hence an isogeny. If $\Theta_J$ is the Abel--Jacobi divisor, then $(u^*\Theta)\Theta_J=2$ and $\Theta_J^2=2$. Hodge index gives $2\deg u=(u^*\Theta)^2\le2$, so $u$ is an isomorphism. The standard Jacobian identities used here hold in characteristic two; see \citet[Theorem~14.23, Corollary~14.24, pp.~230--231; Theorem~14.29, p.~233; Example~14.34, pp.~236--237]{EGM}.
\end{proof}

In the smooth-theta case choose a theta characteristic $\kappa$ on the genus-two curve $G$ and embed it as $x\mapsto\OO_G(x)\otimes\kappa^{-1}$. Such a $\kappa$ exists because multiplication by two on $\Pic^0(G)(k)$ is onto. The resulting $\Theta$ is symmetric. Put $P_a=t_a^*\OO(\Theta)\otimes\OO(-\Theta)$ and write $\Theta_a=\Theta+a$.

For a line bundle $L$, at a point $P$, first jets mean evaluation in $L_P\otimes\OO_{A,P}/\mathfrak m_P^2$, and second jets mean evaluation modulo $\mathfrak m_P^3$. Their dimensions on a surface are three and six, respectively.

\begin{lemma}\label{lem:theta-jets}
In the smooth-theta case the following hold.
\begin{enumerate}
\item Every twist of $2\Theta$ is generated by global sections. The system $|2\Theta|$ separates distinct $x,y$ if $y\ne-x$.
\item The system $|3\Theta|$ separates three distinct values and separates second jets at every point. In particular it is very ample.
\item For distinct $P,Q$, the evaluation on their two first jets is onto except possibly when
\begin{equation}\label{eq:badpairs}
 2P+Q\in\Theta\quad\text{and}\quad P+2Q\in\Theta.
\end{equation}
For a fixed nonzero $g$, only finitely many pairs $(P,P+g)$ satisfy \eqref{eq:badpairs}.
\item If $n\ge4$, $|n\Theta|$ separates any two first jets and any three distinct values.
\end{enumerate}
\end{lemma}
\begin{proof}
The divisors $\Theta_c+\Theta_{-c}$ belong to $|2\Theta|$. To be nonzero at a chosen point, avoid the corresponding two translates of $\Theta$. To vanish at $x$ and not at $y\ne\pm x$, take $c$ on $\Theta+x$ outside $(\Theta+y)\cup(\Theta-y)$. The relevant translates are distinct because the stabilizer of a principal theta divisor is trivial. Their pairwise intersections are finite. Every twist of $2\Theta$ is a translate, since $[2]$ is onto on $k$-points. This proves (1); positive-degree cohomology vanishes by Lemma~\ref{lem:cohomology}.

Two elementary facts about $G$ will also be used. Restriction sends $P_a$ to the degree-zero class $-a$ on $G$, and the Gauss map of $\Theta$ is its canonical map $G\to\PP^1$. For the restriction formula, a generic $a=i(r)-i(s)$ gives the intersection $\Theta\cap\Theta_a=\{i(r),i(\iota s)\}$, where $\iota$ is the hyperelliptic involution. Subtracting the canonical divisor $s+\iota s$ proves $i^*P_{-a}=a$. The difference map $G\times G\to A$ is onto because theta translates intersect, and equality extends from the generic case to all $a$. The Gauss assertion follows by restricting invariant one-forms: $H^0(A,\Omega_A^1)\simeq H^0(G,\omega_G)$. The canonical map is finite and onto in characteristic two as well.

For three distinct points $P_1,P_2,P_3$, choose $a\in\Theta+P_1$ and $b\in\Theta+P_2$, with neither in $\Theta+P_3$ and with $P_3+a+b\notin\Theta$. These are simultaneous nonempty open conditions on the product of the two curves: the addition of two theta translates is onto $A$. The product of the sections of $\Theta_a,\Theta_b,\Theta_{-a-b}$ vanishes at $P_1,P_2$ and not at $P_3$. Permuting the points proves three-value separation.

For first jets at $P$, a theta factor through $P$, multiplied by a section of the remaining twist of $2\Theta$ nonzero there, supplies its Gauss linear form. These linear forms span the cotangent space, and a section nonzero at $P$ supplies the value. The same construction separates two distinct points. To obtain second jets, choose $p,q\in\Theta$, put $a=P-p$, $b=P-q$, and multiply the theta sections with parameters $a,b,-a-b$. The first two vanish at $P$, and the third does not vanish when $3P-p-q\notin\Theta$. This is a dense open of $\Theta^2$. Its image under the product of Gauss maps is dense in $(\PP^1)^2$. Products of the resulting pairs of linear forms span all binary quadratics, also in characteristic two. Hence the quadratic terms, as well as the lower jets already obtained, are arbitrary. This proves (2).

For (3), put $s=P+Q$. Theta translates containing $P,Q$ are parametrized by the length-two intersection $Z=(\Theta+P)\cap(\Theta+Q)$. Its support is preserved by $a\mapsto s-a$. A fixed point has $2a=s$ and corresponds to hyperelliptically conjugate points of $G$ with identical Gauss direction; its intersection multiplicity is at least two. Consequently two distinct support points are exchanged, and a unique support point is fixed.

For $a$ in this support, restriction to $\Theta_a$ gives a surjection
\[
 H^0(3\Theta)\longrightarrow H^0\bigl(G,\omega_G^3\otimes(-3a)\bigr).
\]
Write $p=P-a$, $q=Q-a$ on $G$. The degree-six bundle has vanishing $H^1$, and the cokernel of its four tangential jet evaluations is the $H^1$ of the residual degree-two bundle. This has dimension one precisely when that bundle is canonical. Thus the tangential jet conditions have rank four unless
$\omega_G^3\otimes(-3a)(-2p-2q)=\omega_G$, when the rank is three. Since the class of $p+q-K_G$ is $s-2a$, the exception is exactly $a=2s$. The sections containing $\Theta_a$ supply the two normal derivatives: their residual bundle is $2\Theta\otimes P_a$. By (1), these two derivatives are independent unless $a=-s$. Thus the full six jets are obtained whenever a support point avoids both $-s$ and $2s$.

If every support point lies among these two bad parameters, then $-s$ is a support point. Indeed, if the sole support point is $2s$, its fixed-point equation gives $4s=s$ and hence $2s=-s$. Therefore rank failure implies that $\Theta_{-s}$ contains both $P,Q$, which is \eqref{eq:badpairs}. For $Q=P+g$, these conditions are $3P+g\in\Theta$ and $3P+2g\in\Theta$. The two translates are distinct when $g\ne0$, and $[3]$ is finite; their common inverse image is finite.

Finally let $n=3+r$ with $r\ge1$. Choose $(a,b)$ satisfying $3b=-ra$. The parameter space is finite and surjective over either factor, and each component dominates both. The residual factorization is
\[
 \OO(n\Theta)=t_a^*\OO(r\Theta)\otimes t_b^*\OO(3\Theta).
\]
The distinguished section of the first factor is nonzero at prescribed finitely many points on an open in $a$. For a prescribed pair $P,Q$, the bad translated cubic parameters satisfy the two conditions \eqref{eq:badpairs} with $3b$ added; they form a finite set of $b$. Choose a pair $(a,b)$ outside these bad loci. Multiplication by the nonvanishing residual section transfers all six jets. The identical argument transfers the three independent values from (2). This proves (4).
\end{proof}

On a product $E_1\times E_2$, $\OO(3\Theta)=\OO_{E_1}(3O_1)\boxtimes\OO_{E_2}(3O_2)$ is very ample by the plane cubic embeddings. Its Frobenius fibres have length four and are products of length-two elliptic fibres. A degree-three line bundle on an elliptic curve evaluates onto every length-two subscheme, since the kernel has degree one and vanishing $H^1$. We record the consequence also in the smooth-theta case.

\begin{corollary}\label{cor:fgenerated}
For relative Frobenius $F:A\to A^{(2)}$ on a principally polarized surface, $F_*\OO_A(3\Theta)$ is generated by its global sections.
\end{corollary}
\begin{proof}
The local algebra of a Frobenius fibre is $k[x,y]/(x^2,y^2)$, a quotient of $k[x,y]/(x,y)^3$. Lemma~\ref{lem:theta-jets} supplies the required evaluation in the smooth-theta case. The preceding product calculation supplies it in the other case. Nakayama's lemma then proves global generation.
\end{proof}

\section{The Cartier tensor and a relative determinant}

Put $L=\OO_A(3\Theta)$ and $W=H^0(A,L)$, of dimension nine. Fix invariant volume forms on $A$ and $A^{(2)}$. Relative Frobenius duality and $F^*L^{(2)}=L^2$ give a pairing
\begin{equation}\label{eq:beta}
 \beta:W\otimes W\longrightarrow W^{(2)},\qquad
 (s,t)\longmapsto\Tr_F(st).
\end{equation}
The trace here is the trace on top forms, transported by the chosen trivializations, not the field trace of a purely inseparable extension. Locally it is the coefficient-of-$xy$ map on a two-element $2$-basis. In particular the sheaf pairing
\begin{equation}\label{eq:perfect}
 F_*L\otimes F_*L\longrightarrow L^{(2)}
\end{equation}
is perfect: in the basis $1,x,y,xy$, each monomial pairs nontrivially with its complementary monomial. This is the explicit finite-Frobenius duality of \citet[Lemma~1.3.6, Proposition~1.3.7 and Theorem~1.3.8, pp.~26--28]{BrionKumar}.

\begin{lemma}[The typed symmetric tensor]\label{lem:tensor}
There is a basis $(e_a)_{a\in\FF_3^2}$ of $W$ identifying $W^{(2)}$ with $W^*$ under the finite level-three Heisenberg group such that the trilinear tensor associated to \eqref{eq:beta} is fully symmetric. It satisfies
\[
 T(e_a,e_a,e_b)=c\,\delta_{ab},
\]
where $c=0$ exactly when $A$ is nonordinary. In the nonordinary case $T$ is alternating and is specified by four scalars $\alpha_i$, one for each direction $i\in\PP^1(\FF_3)$: its value on three distinct basis vectors is $\alpha_i$ when their indices lie on a line of direction $i$, and is zero otherwise.
\end{lemma}
\begin{proof}
The theta group of $L$ has quotient $A[3]$ and admits a theta structure. Its weight-one representation on $W$ is the Schr\"odinger representation, by \citet[Lemma~8.24, p.~121; Theorem~8.28 and description~8.29, pp.~122--123; Corollary~9.17, pp.~133--134]{EGM}. Fix $\zeta\in\mu_3(k)\setminus\{1\}$. The finite subgroup generated by
\[
 e_a\longmapsto e_{a+x},\qquad
 e_a\longmapsto\zeta^{y\cdot a}e_a,\qquad
 e_a\longmapsto\zeta e_a
\]
has order $3^5$. If $\rho$ denotes its matrices, the Frobenius-twisted matrices obey $\rho(h)^{[2]}=\rho(h)^{-t}$. This holds directly for each displayed generator, since $\zeta^2=\zeta^{-1}$ and the translation matrices are orthogonal. Thus the indicated basis pairing identifies $W^{(2)}$ with the contragredient representation. This statement is restricted to the finite central $\mu_3$; it is not an assertion about the whole scalar $\mathbf G_m$.

Functoriality of the Frobenius trace and projection formula make \eqref{eq:beta} equivariant for this twisted action. Accordingly $T$ is invariant under the finite group and symmetric in its first two arguments. Modulations force a nonzero coefficient to have indices $a+b+c=0$. Translation by $b-a$ sends such a triple to $(b,c,a)$. Cyclic symmetry, together with the first transposition, gives full symmetry. If two indices coincide, the sum-zero condition forces all three to coincide. Projection formula gives $\beta(s,s)=c s^{(2)}$, where $c$ is the coefficient of Cartier on the invariant top form. Hence the diagonal coefficients are as asserted. This scalar is nonzero exactly when $A$ is Frobenius split: by finite Frobenius duality and $\omega_A\simeq\OO_A$, the space of maps $F_*\OO_A\to\OO_{A^{(2)}}$ is one-dimensional, and evaluation at $1$ is multiplication by $c$. A nonzero evaluation can be rescaled to split the unit map. Frobenius splitting is equivalent to ordinarity for an abelian variety \citep[Definition--Proposition~2.1, p.~1071]{SannaiTanaka}. When it is zero the tensor is alternating. Every triple of distinct indices with sum zero is an affine line, and translations identify coefficients on parallel lines, giving the four parameters.
\end{proof}

The following cubic is the reduction modulo two of the known Coble-cubic formula of \citet[equation~(3.9), author version p.~14]{GrusonSam}. We give its polynomial proof and characteristic-two exceptional-tensor analysis in full. For $I=\FF_3^2$, let
\[
 P_*(s)=\sum_{a\in I}s_a^3,\qquad
 P_i(s)=\sum_{\ell\parallel i}\prod_{a\in\ell}s_a.
\]
The second sum has three terms, one for each affine line in direction $i$.

\begin{lemma}[Pfaffian identity and its exceptional tensors]\label{lem:pfaffian}
Let $T$ be the alternating tensor with the four parameters of Lemma~\ref{lem:tensor}, and set $M(s)_{ab}=T(s,e_b,e_a)$. Then
\begin{equation}\label{eq:q}
 q_\alpha(s)=\left(\prod_i\alpha_i\right)P_*(s)
 +\sum_i\alpha_i\left(\sum_{j\ne i}\alpha_j^3\right)P_i(s)
\end{equation}
satisfies
\[
 M(s)s=0,\qquad
 \operatorname{Pf}(M(s)_{\widehat a,\widehat a})=q_\alpha(s)s_a,
 \qquad\operatorname{adj}M(s)=q_\alpha(s)^2ss^t.
\]
For $\alpha\ne0$, $q_\alpha$ is identically zero precisely at forty projective parameter points: the four points with one nonzero coordinate, and the thirty-six points with three nonzero coordinates whose cubes are equal. At every such point, $W$ is an orthogonal direct sum of three three-dimensional spaces on each of which $T$ is a nonzero volume form.
\end{lemma}
\begin{proof}
Perform the calculation over $\FF_2[\alpha_i,s_a]$. Alternation gives $Ms=0$. The vector $v$ of principal eight-dimensional Pfaffians is also in the kernel, by the Pfaffian expansion identity. At $s=e_0$, the minor omitting zero consists of four alternating two-by-two blocks, one for each direction; its Pfaffian is $\prod_i\alpha_i$. Thus the universal matrix has rank eight, so $v=qs$ in the fraction field. The polynomial identities $v_as_b=v_bs_a$ and the pairwise coprimality of the variables $s_a$ imply that $s_a$ divides $v_a$. Therefore $q$ is a polynomial, of degree three in $s$ and four in $\alpha$.

The finite Heisenberg generators preserve $T$ and have determinant one. The transformation rule for $\bigwedge^8W^*\simeq W\otimes\det W^*$ shows that $v$ transforms as $s$, and hence $q$ is invariant. A modulation-invariant cubic monomial has its three indices summing to zero. It is either $s_a^3$ or the product on an affine line. Translation invariance therefore puts $q$ in the span of $P_*,P_i$. The calculation at $e_0$ gives its $P_*$ coefficient.

For a line coefficient, set all section coordinates outside the vertical line $x=0$ equal to zero, and denote its three remaining coordinates by $x_0,x_1,x_2$. The minor omitting $(0,0)$ has a two-by-two block of entry $\alpha_\infty x_0$ and an off-diagonal six-by-six block whose Pfaffian is the determinant of
\[
 B=\begin{pmatrix}
 \alpha_0x_0&\alpha_1x_2&\alpha_2x_1\\
 \alpha_2x_2&\alpha_0x_1&\alpha_1x_0\\
 \alpha_1x_1&\alpha_2x_0&\alpha_0x_2
 \end{pmatrix}.
\]
Expansion of its six terms gives
\[
 \det B=\alpha_0\alpha_1\alpha_2(x_0^3+x_1^3+x_2^3)
 +(\alpha_0^3+\alpha_1^3+\alpha_2^3)x_0x_1x_2.
\]
The principal Pfaffian is $\alpha_\infty x_0\det B$; dividing that Pfaffian by $x_0$ gives $q|_{x=0}=\alpha_\infty\det B$, and hence the coefficient of $P_\infty$ in \eqref{eq:q}. Linear permutations of the index plane give the other three coefficients. The identity $\operatorname{adj}M=vv^t$ for odd alternating matrices then gives the adjugate formula.

The five orbit sums have disjoint supports. Consequently $q=0$ is equivalent to
\[
 \prod_i\alpha_i=0,\qquad
 \alpha_i\sum_{j\ne i}\alpha_j^3=0\quad\text{for each }i.
\]
If $r$ coordinates are nonzero, these equations make all their cubes equal to a common $u\ne0$ and require $ru=u$ in characteristic two. Thus $r$ is odd, and the product condition leaves $r=1$ or $3$. The projective counts are $4$ and $4\cdot3^2$.

When $r=1$, the three parallel lines partition the basis and give the required decomposition. When $r=3$, rescale so that all nonzero coefficients lie in $\mu_3$. A basis change $e_a\mapsto\zeta^{Q(a)}e_a$ with $Q$ a binary quadratic form over $\FF_3$ multiplies the coefficient in direction $v$ by $\zeta^{-Q(v)}$, because $\sum_{t\in\FF_3}Q(a+tv)=-Q(v)$. Evaluation of binary quadratics at three distinct directions is an isomorphism, so all three coefficients can be made one. Move the missing direction to the vertical direction by a linear change of the index plane.

For the final basis change, start with the sum of the three volume forms on horizontal lines and set
\[
 f_{x,u}=\sum_{y\in\FF_3}\zeta^{uy}e_{x,y}.
\]
This Fourier matrix is invertible; its inverse uses $\zeta^{-uy}$, since $3=1$ in $k$. Its transformed tensor is one exactly when the three first indices are distinct and the second indices sum to zero, and zero otherwise. These are precisely the lines of the three nonvertical directions. Thus the normalized three-direction tensor is equivalent to the one-direction tensor. Transporting its three orthogonal summands proves the last assertion.
\end{proof}

\begin{proposition}[The cubic principal seed]\label{prop:seed}
For every principally polarized abelian surface, a generic smooth curve in $|3\Theta|$ has genus ten and $2$-rank $8+f_A$. Its relative Frobenius is bijective on its eight-dimensional quotient.
\end{proposition}
\begin{proof}
The pairing \eqref{eq:perfect} and Corollary~\ref{cor:fgenerated} imply that $\beta$ is nonzero. Suppose first that $A$ is nonordinary. If $q_\alpha=0$, Lemma~\ref{lem:pfaffian} gives a rank-three projector $D$ on $W$ such that
\[
 T(Ds,t,u)=T(s,Dt,u)
\]
and whose summands are mutually orthogonal. It preserves the kernel of every fibre evaluation $W\to(F_*L)_x$. Indeed, if $a+b$ vanishes in the fibre with $a$ in one summand and $b$ in its orthogonal complement, the value of $a$ pairs to zero with the complement by orthogonality and with its own summand because $a+b$ has zero value. Global generation and perfectness force $a$ to have zero value. Thus the kernel bundle of $W\otimes\OO\twoheadrightarrow F_*L$ is preserved by $D$: the induced map from it to $F_*L$ vanishes at every closed point and hence identically. The projector descends to a nontrivial endomorphism of $F_*L$.

Oda's theorem contradicts this. The restriction of $\lambda_L=3\lambda_\Theta$ to the scheme-theoretic kernel of $F$ is a closed immersion, because $\lambda_\Theta$ is an isomorphism and multiplication by three is an automorphism on $A[2]$. Hence $\End(F_*L)=k$ by \citet[Theorem~1.2(i), p.~45; proof pp.~43--45]{Oda}. A scalar cannot act on $W$ as a rank-three projector. Therefore $q_\alpha\ne0$.

It remains essential to compute the map on the relative cohomology. By Lemma~\ref{lem:cohomology} and Serre duality, the fibre is the hyperplane $\ann(s)\subset W^*$, and its Frobenius-domain hyperplane is $\ann(s^{[2]})$. The contraction matrix $M(s)$ is precisely the linearization of Frobenius: duality converts the multiplication by the defining section in the restriction sequence to \eqref{eq:beta}. Since $Ms=0$, it maps the indicated domain into the indicated target. On $s_0=1$, the domain coordinate basis is $e_j^*+s_j^2e_0^*$ and the target basis is $e_i^*+s_ie_0^*$, for $1\le i,j\le8$. The relative matrix is
\[
 R(s)_{ij}=M(s)_{ij}+s_j^2M(s)_{i0}.
\]
Expansion by the rank-one column additions, using the adjugate in Lemma~\ref{lem:pfaffian}, gives
\begin{equation}\label{eq:relative-det}
 \det R(s)=q_\alpha(s)^2\sum_{a\in I}s_a^3.
\end{equation}
For completeness, the unmodified minor contributes $q^2s_0^2=q^2$. Replacing its $j$th column by $s_j^2(M_{i0})_{i\ne0}$ contributes $s_j^2q^2s_0s_j=q^2s_j^3$; terms with two replaced columns vanish. This proves \eqref{eq:relative-det}. Both factors are nonzero polynomials. Thus the relative map is generically invertible, which proves stable rank eight rather than only an initial rank assertion.

If $A$ is ordinary, Lemma~\ref{lem:tensor} gives a symmetric contraction matrix with diagonal $c s_a$, where $c\ne0$. In characteristic two, terms of a symmetric determinant with a permutation cycle of length at least three cancel with their inverse-cycle partners. The remaining terms are products of diagonal entries and squares of off-diagonal entries. The monomial $\prod_as_a$ therefore has coefficient $c^9\ne0$, since every nondiagonal term has some positive even exponent. The full matrix is generically invertible. It preserves the relative hyperplanes: full symmetry gives $M(s)s=c s^{[2]}$. Therefore its restriction between these equal-dimensional hyperplanes is also invertible. This proves the same relative assertion.

Finally $|3\Theta|$ is very ample, so its smooth-member open meets the nonvanishing determinant open. Equation~\eqref{eq:dimensions} gives genus ten and $2$-rank $8+f_A$.
\end{proof}

\section{Finite norms and their singularities}

\begin{lemma}[Ordinary principal members]\label{lem:ordinary-principal}
If $(B,\Theta)$ is an ordinary principally polarized surface, a generic member of $|n\Theta|$ is smooth and ordinary for every $n\ge3$. The line bundle $\OO(n\Theta)$ is very ample.
\end{lemma}
\begin{proof}
In the smooth-theta case $\Theta$ is a genus-two curve whose Jacobian is $B$, so it is ordinary. Choose $n$ translates $\Theta_{a_i}$ with $\sum_i a_i=0$. Their union lies in $|n\Theta|$. The parameters can be chosen so that all pairwise intersections are transverse and no triple intersection occurs. Pairwise transversality holds on a dense open of the difference parameter: the difference map $\Theta\times\Theta\to B$ has nonzero differential where the Gauss directions differ, and the Gauss map is onto. For $n\ge4$, the sum-zero parameter space projects onto any three coordinates, so the usual pair and triple open conditions can be imposed simultaneously.

For $n=3$, a pairwise difference is still freely variable, so the transversality condition is proper. A triple intersection with common point $P$ gives points $p_i\in\Theta$ with $a_i=P-p_i$ and $3P=p_1+p_2+p_3$. This incidence has dimension three, since $[3]$ is finite, whereas the sum-zero parameter space has dimension four. Thus triples can also be avoided. Lemma~\ref{lem:nodal} makes the resulting union ordinary. In the product case take $n$ distinct vertical and $n$ distinct horizontal ordinary elliptic curves; choose their coordinate sums so that the union belongs to $|n\Theta|$. This is an ordinary nodal grid.

For $n=3$, very ampleness was proved above. For $n>3$, factor $\OO(n\Theta)$ as a translated $3\Theta$ bundle times a translated $(n-3)\Theta$ bundle, with parameters satisfying $3b=-(n-3)a$. The second factor has a section nonzero at any prescribed finite set of points. Multiplication transfers point and tangent separation from the first factor. Thus $n\Theta$ is very ample. Openness from Lemma~\ref{lem:stable} now smooths the ordinary nodal member to a smooth ordinary member.
\end{proof}

\begin{lemma}[Nodal norm images]\label{lem:norm-nodes}
Let $\phi:B\to A$ be a finite \'etale isogeny, where $(B,\Theta)$ is principally polarized. For every $n\ge3$, a general smooth $C\in|n\Theta|$ maps birationally to an integral nodal curve, with exactly $n^2(\deg\phi-1)$ nodes. This assertion allows even $\deg\phi$.
\end{lemma}
\begin{proof}
Let $K=\ker\phi(k)$, a finite constant group. We exclude tangencies and triple identifications for the finitely many nonzero translations in $K$.

Suppose first that $\Theta$ is smooth. For $n=3$, Lemma~\ref{lem:theta-jets} shows that the six first-jet coordinates at $P,P+g$ are independent outside a finite set of $P$. For $n\ge4$ they are independent everywhere. At such a pair, requiring the two values to vanish gives two independent equations, and requiring proportional gradients gives a further nonzero determinant equation. Thus the tangency incidence has codimension three in the section space while $P$ varies in dimension two. At the finite exceptional set for $n=3$, the two independent value equations suffice to avoid the pair altogether. For distinct $0,g,h\in K$, three-value separation imposes three conditions at $P,P+g,P+h$ while $P$ varies in dimension two. Hence triple identifications are avoided on a nonempty open.

We spell out the product case, where cubic elliptic factors have a smaller jet space. A degree-$n$ line bundle on an elliptic curve separates two values and separates a first jet at a single point for $n\ge3$. For $n=3$, the evaluation on the two first jets at distinct $x,y$ has rank three; after their two values are prescribed to be zero, the two derivatives form a line. Both derivative coefficients on this line are nonzero unless the unique residual zero equals $x$ or $y$. For $y=x+g$, these exceptions are the finite equations $3x+g=0$ or $3x+2g=0$, after the chosen origin is incorporated in the line bundle class.

When $g=(g_1,g_2)$ has both coordinates nonzero, at points avoiding those finite exceptions in both factors the tensor-product section space supplies the full six surface first jets. To see this, split each three-dimensional factor space into its two value coordinates and its one-dimensional common derivative direction. The tensor of the two value spaces supplies both surface values; derivative-direction times value-space supplies independent derivatives at the two points in one direction, and value-space times derivative-direction does the same in the other direction. It also suffices that just one factor have both derivative coefficients nonzero. To verify this, use value bases $u_1,u_2$ and $v_1,v_2$ and residual derivative directions $u_3,v_3$. Suppose the second factor has both coefficients nonzero. Modulo the two surface values, the two first-direction derivatives involve respectively the disjoint pairs of coordinates $(u_3v_1,u_2v_1)$ and $(u_3v_2,u_1v_2)$, with each pair having a nonzero coefficient because each individual elliptic first jet is onto. The second-direction derivatives can then be adjusted independently by $u_1v_3$ and $u_2v_3$. All four derivatives are therefore independent. The locus where both factors fail is finite and is avoided by the two independent value conditions.

If precisely one $g_i$ is nonzero, the cubic tensor evaluation has rank five. After the two values vanish, the two derivatives in the varying coordinate lie on a nonzero line, and the two transverse derivatives are independent. Their tangency determinant is a nonzero polynomial in these three parameters. This again gives three conditions. The possible further derivative exceptions occur on a one-dimensional locus of $P$; there the same nonzero derivative line and independent transverse derivatives still give the third condition. Thus its incidence is proper as well.

For triples in a cubic product, three Segre value tensors are dependent only if one coordinate is constant and the three value vectors in the other plane cubic are collinear. Indeed, if two vectors in each factor are independent, their two simple tensors span no third simple tensor except on the two ruling lines; this follows immediately in the resulting two-by-two tensor matrix. A plane cubic triple $x,x+g,x+h$ is collinear only when $3x+g+h=0$. Hence the rank-two exceptional locus has dimension at most one and the two value conditions exclude it generically. Elsewhere the triple values have rank three and the point parameter has dimension two.

For $n\ge4$ on a product, three distinct elliptic values are independent. For $n\ge5$, two elliptic first jets are independent as well. For $n=4$, their rank can drop from four to three only if the line bundle is $\OO(2x+2y)$; for $y=x+g$ this is a finite condition $4x+2g=0$. Tensor products therefore give six surface first jets away from a finite set when both $g_i$ are nonzero. When one $g_i$ vanishes, the exceptional locus has dimension one and rank five, with the same nonzero tangency determinant described above. Three distinct surface values are independent: separate the first-coordinate classes and then the at most three second-coordinate values within each class. These assertions follow in each elliptic factor from Riemann--Roch applied to the evaluation kernel.

All the undesirable incidences are thus proper. A general section also has no nontrivial kernel stabilizer. If translation by $g$ does not preserve the line bundle it preserves no member; if it does, its projective action in the very ample embedding is nontrivial, and its invariant members form a proper closed subset. Therefore $C\to\phi(C)$ has generic degree one. Since $\phi$ is \'etale, the only singularities of the image are the transverse pairs just obtained, so they are nodes. For each nonzero $g$, the curves $C$ and $C+g$ meet in $2n^2$ points. These are ordered branch pairs; every node contributes two. Consequently the number of nodes is $n^2(|K|-1)$. This is a count of finite sets and does not divide by two in the coefficient field.
\end{proof}

\begin{lemma}[Principal norm sources]\label{lem:principal-source}
Let $M$ be ample on $A$ and put $D=\chi(M)$. Suppose either that $D$ is odd or that $A$ is ordinary. There are a principally polarized surface $(B,\Theta_B)$ and a finite \'etale isogeny $\phi:B\to A$ of degree $D$ such that
\[
 \phi^*M\equiv\OO_B(D\Theta_B),\qquad
 \Nm_\phi\OO_B(\Theta_B)\equiv M.
\]
There is also an ample bundle $N$ on $A^\vee$ with
\begin{equation}\label{eq:dual-pol}
 \chi(N)=D,\qquad\lambda_N=D\lambda_M^{-1}.
\end{equation}
\end{lemma}
\begin{proof}
Write $K=\ker\lambda_M$. If $D$ is odd, the finite symplectic group $K$ has type $(a,b)$ with $a\mid b$, $ab=D$, and is killed by $D$. If $A$ is ordinary, its $2$-power torsion over $k$ is the product of a multiplicative and a constant \'etale factor. The $2$-primary part of $K$ decomposes accordingly. The perfect alternating commutator pairing is trivial on each factor and pairs them perfectly with each other. Each therefore has order equal to the $2$-part $D_2$ of $D$, and is killed by $D_2$. On the odd-primary components the usual symplectic decomposition gives exponent dividing the corresponding part of $D$. Thus $D$ kills $K$ in both cases.

The commutator pairing of $M^D$ on $K$ is the $D$th power of that of $M$ and is trivial. Finite isotropic descent therefore descends $M^D$ through $\lambda_M:A\to A^\vee$ to a line bundle $N$; the exact descent statement, including nonreduced kernels over an algebraically closed field, is \citet[Theorem~8.10 and Corollary~8.11, pp.~116--117]{EGM}. Since the pullback is ample, $N$ is ample. The pullback polarization identity gives
$\lambda_M\lambda_N\lambda_M=D\lambda_M$, and cancellation gives \eqref{eq:dual-pol}. Euler characteristics give $\chi(N)=D^3/D^2=D$.

Choose a maximal isotropic subgroup $H\subset\ker\lambda_N$. When $D$ is odd every such $H$ is \'etale of order $D$. When $A$ is ordinary, choose its $2$-primary factor to be the full connected multiplicative factor, and choose a maximal isotropic in each odd-primary factor. This $H$ has order $D$ and has \'etale Cartier dual. The order and isotropic assertions follow from the perfect pairing; see also \citet[Proposition~8.19--Corollary~8.20, p.~119]{EGM}.

Let $\psi:A^\vee\to B^\vee=A^\vee/H$ be the quotient. Isotropic descent gives a principal polarization $\Theta_{B^\vee}$ with $\psi^*\Theta_{B^\vee}\equiv N$, since its Euler characteristic is $D/|H|=1$. Give $B$ the dual principal polarization and let $\phi=\psi^\vee$. Isogeny duality identifies $\ker\phi$ with $H^D$, so $\phi$ is \'etale of degree $D$; see \citet[Theorem~7.5, p.~100]{EGM}. From
$\lambda_N=\psi^\vee\lambda_{\Theta_{B^\vee}}\psi$
and \eqref{eq:dual-pol}, cancellation yields
$\psi\lambda_M\psi^\vee=D\lambda_{\Theta_B}$.
Hence $\phi^*M\equiv D\Theta_B$. Pulling back the norm of $\Theta_B$ gives the sum of its $D$ translates, numerically $D\Theta_B$. Since pullback on the torsion-free N\'eron--Severi group is injective, its norm is numerically $M$.
\end{proof}

\begin{lemma}\label{lem:exact-system}
The norm of a very ample line bundle under a finite \'etale morphism of projective varieties over $k$ is very ample. On an abelian variety, any effective divisor in the numerical class of an ample $N$ can be translated into the exact system $|N|$.
\end{lemma}
\begin{proof}
To separate two target points by a norm section, choose a source section vanishing at one selected point over the first and nonzero at all the other finitely many points over both targets. Very ampleness makes each unwanted vanishing a proper hyperplane in the space of sections vanishing at the selected point. To separate a tangent vector, choose its lifted derivative nonzero at the selected point and keep all other fibre values nonzero. The norm is locally a product; its other factors are units, so its derivative is the selected nonzero derivative times a unit. Basepoint freeness is proved by a section avoiding the entire finite fibre. Thus norm sections themselves separate points and tangent vectors.

For the second assertion, the two line bundles differ by an element of $\Pic^0$. Surjectivity of $\lambda_N$ realizes that difference by a translation of $N$, and translating the divisor gives the required member. Translations preserve smoothness, nodes and cohomological ranks.
\end{proof}

\section{The integral even-degree bound and the proofs}

\begin{lemma}[Connected Prym bound]\label{lem:prym}
Let $M$ be ample of odd Euler characteristic $D$ and let $C\in|M^m|$ be smooth. If $m$ is even, then
\[
 g(C)-f_C\ge2(2-f_A).
\]
For any smooth ample curve on $A$, one always has $g(C)-f_C\ge2-f_A$.
\end{lemma}
\begin{proof}
Restriction $j:A^\vee\to J_C$ is a closed immersion. For a nontrivial $P\in\Pic^0(A)(k)$, $H^0(A,P)=0$ and $H^1(A,P(-C))=0$ by ample vanishing and duality, so $H^0(C,P|_C)=0$. Thus the geometric point kernel of $j$ is trivial. Its tangent map is injective because $H^1(A,\OO_A(-C))=0$. Its finite kernel is therefore \'etale with only the identity point, and hence trivial.

Let $q:J_C\to A$ be its dual under the principal polarization of $J_C$. It is onto with smooth connected kernel $P_C$. To justify the last assertion, quotient $J_C$ by the reduced identity component of $\ker q$. The remaining map to $A$ is an isogeny; its dual kernel embeds in $\ker j=0$, so isogeny duality makes it an isomorphism. Consequently $\dim P_C=g(C)-2$ and $f_C=f_{P_C}+f_A$. The latter follows also from Poincar\'e reducibility and isogeny invariance of $2$-rank, for which see \citet[Proposition~5.22, pp.~80--81]{EGM}. This already proves the unconditional bound.

Let $N$ be the integral dual polarization of Lemma~\ref{lem:principal-source}. The exact induced-polarization formula is
\begin{equation}\label{eq:prym-mult}
 qj=j^*\lambda_{J_C}=m\lambda_N.
\end{equation}
Indeed Debarre's norm endomorphism formula in dimension two gives $\alpha(C,M)=mD\,\mathrm{id}_A$ and
$\lambda_M(j^*\lambda_{J_C})\lambda_M=mD\lambda_M$
\citep[\S3, equations~(3.2)--(3.3), pp.~346--347]{Debarre}. Cancelling in the rational homomorphism group gives $j^*\lambda_{J_C}=mD\lambda_M^{-1}=m\lambda_N$. Both sides are integral homomorphisms, so \eqref{eq:prym-mult} is an integral identity, including when $m$ is even.

For even $m$, this identity embeds $A^\vee[2]$ in $P_C[2]$. Apply Frobenius twist to that immersion and use functoriality of Verschiebung. The identity component of $\ker V_{A^\vee}$ embeds in that of $\ker V_{P_C}$. For an abelian variety $B$, this connected kernel has order $2^{\dim B-f_B}$: $V_B$ has degree $2^{\dim B}$ and its \'etale degree is $2^{f_B}$, as follows from the Frobenius--Verschiebung relations and the torsion description in \citet[Propositions~5.20--5.22, pp.~80--81]{EGM}. Hence $\dim P_C-f_{P_C}\ge2-f_A$. Adding the ambient defect proves the even-degree bound.
\end{proof}

\begin{proof}[Proof of Theorem~\ref{thm:odd}]
First suppose $d$ is odd. Apply Lemma~\ref{lem:principal-source} to $M^d$, whose Euler characteristic is $Dd^2$, to obtain an \'etale norm source of degree $Dd^2$. Proposition~\ref{prop:seed} supplies an open of smooth curves in $|3\Theta_B|$ with defect $2-f_B=2-f_A$. Lemma~\ref{lem:norm-nodes} supplies another nonempty open whose norm images are integral nodal. A curve in their intersection has a norm image with $9(Dd^2-1)$ nodes. Its arithmetic genus is $10+9(Dd^2-1)=9Dd^2+1$, and its defect is still $2-f_A$ by Lemma~\ref{lem:nodal}.

Its numerical class is $M^{3d}$, and Lemma~\ref{lem:exact-system} translates it into that exact system. The same lemma proves $M^{3d}$ very ample by taking norms of $3\Theta_B$. Openness therefore gives smooth members of defect at most $2-f_A$. The unconditional Prym bound forces equality. This proves the odd-degree formula in every odd degree, including one.

Now suppose $d$ is even. Choose a smooth curve of the just-proved minimal defect in $|M^3|$ and another in $|M^{3(d-1)}|$. They may be chosen to meet transversely, since these systems are very ample and the desired cohomological condition is open. Their union is a connected nodal member of $|M^{3d}|$ with defect $2(2-f_A)$. Very ampleness of $M^3$ implies that of its positive powers, so openness gives a smooth member of no greater defect. Lemma~\ref{lem:prym}, with $m=3d$, gives the reverse inequality. Thus equality holds generically. Finally subtract $f_A$ from the resulting full $2$-rank and use \eqref{eq:dimensions} to obtain the displayed relative formulas.
\end{proof}

\begin{proof}[Proof of Theorem~\ref{thm:ordinary}]
Use the ordinary case of Lemma~\ref{lem:principal-source} for the given ample $M$. Its principal source $B$ is ordinary, since $2$-rank is invariant under isogeny. For $n\ge3$, intersect the ordinary smooth-member open of Lemma~\ref{lem:ordinary-principal} with the nodal-norm open of Lemma~\ref{lem:norm-nodes}. The image of such a curve has ordinary normalization and only nodes, so is ordinary as a nodal curve by Lemma~\ref{lem:nodal}. Its genus is
\[
 (n^2+1)+n^2(D-1)=Dn^2+1.
\]
Its numerical class is $M^n$. Lemma~\ref{lem:exact-system} translates it into the exact linear system and proves that $M^n$ is very ample. Openness then gives a smooth ordinary member of that system and hence the generic assertion. Equation~\eqref{eq:dimensions} gives the relative rank.
\end{proof}

\begin{proof}[Proof of Corollary~\ref{cor:parity}]
For $D=1$ the relative dimension is $9d^2-1$. The difference from the stable rank is zero in odd degree and $2-f_A>0$ in even degree. The formula is positive already at $d=1$ and increases thereafter.
\end{proof}

\section{Ambient hypersurfaces and scope}

\begin{proposition}\label{prop:ambient}
Suppose $A\hookrightarrow\PP^N$ is the embedding defined by a very ample $L$. If
\[
 H^0(\PP^N,\OO(d))\longrightarrow H^0(A,L^d)
\]
is onto, the generic rank in the complete system equals the generic rank obtained from ambient degree-$d$ hypersurfaces. A sufficient condition is
$d\ge\max\{1,\operatorname{reg}(\mathcal I_A)-1\}$.
For the cubic Segre embedding of a product of elliptic curves the restriction map is onto in every positive degree.
\end{proposition}
\begin{proof}
A surjective linear restriction map has a dominant projective map on the open where restriction is nonzero, with irreducible fibres. The inverse image of a nonempty open of the complete system is therefore nonempty open. The regularity bound gives $H^1(\mathcal I_A(d))=0$ and hence surjectivity by the ideal sequence.

For an elliptic plane cubic, $0\to\OO_{\PP^2}(d-3)\to\OO_{\PP^2}(d)\to\OO_E(3dO)\to0$ and $H^1(\PP^2,\OO(d-3))=0$ give surjectivity in every degree. The degree-$d$ Segre coordinates span all products of degree-$d$ coordinates in the two factors. Tensoring the two elliptic surjections gives the assertion for the cubic Segre surface.
\end{proof}

Applying this proposition with $L=M^3$ gives the ambient-hypersurface version of Theorem~\ref{thm:odd} in the indicated range. If an ample $M$ itself gives the chosen embedding of an ordinary surface, Theorem~\ref{thm:ordinary} transfers for
$n\ge\max\{3,\operatorname{reg}(\mathcal I_A)-1\}$.
No ambient restriction-surjectivity assertion is implicit outside these ranges.

The proofs do not determine the generic rank for an arbitrary even-degree polarization on a nonordinary surface, nor for every primitive low-degree linear system. They also do not determine a monodromy group. The determinant in \eqref{eq:relative-det} is used to prove invertibility on relative cohomology; knowledge of an initial matrix rank alone would not supply its stable rank. These distinctions delimit the stated conclusions.

\section*{Sources, computation, and method}

The geometric inputs from the literature are the theta representation and finite isotropic descent in \citet{EGM}, the integral norm-polarization calculation in \citet{Debarre}, Frobenius duality in \citet{BrionKumar}, and Oda's simplicity theorem \citep{Oda}. Their exact statements and the relevant proofs were checked at the cited locations. The known Gruson--Sam Pfaffian identity has been explicitly credited and rederived over characteristic two. The characteristic-two jet arguments, the exceptional tensor decomposition, the relative determinant, the nodal incidence arguments, and the rank deductions have been written out here. The original manuscript passage was checked against the Montpellier scan. A bounded search of journals, author pages, arXiv, Hexagon and ECCC found no exact match for the two displayed all-degree rank theorems. Searches of recent AI mathematics, including OpenAI publications, identified no directly relevant theorem about the generic stable ranks in these linear systems. No novelty or priority is inferred from those search results.

The polynomial calculation is proved without computation. Reproducible exact arithmetic accompanying the paper checks the nine principal Pfaffians, the forty exceptional projective tensors, the Fourier change of basis, and the displayed genus and rank identities. Such finite checks supplement the identities and do not prove the geometric existence or all-degree statements. Source code and records are available through \url{\workbenchurl} and \url{\recordurl}.

The research and writing were performed by GPT-6 Astra Pro in ChatGPT Pro mode. The complete proofs and primary sources received targeted same-model checks. Separate proof-unseen statement-only derivations established every main target: the ordinary arbitrary-ample theorem, the integral even upper bound, principal cubic attainment, the norm-image construction for every odd Euler characteristic, all parity cases, and the ambient restriction qualification. Their complete proofs were read and compared, with no defect found in those stated scopes. The separate derivations were frozen before comparison. Reasoning-effort metadata is not exposed. Earlier records attributed targeted checks to Claude Opus~5.5 at effort max; those historical records are distinct from the current checks. No human has checked the mathematics. This work is dedicated to the public domain under CC0~1.0.

\begin{thebibliography}{9}
\bibitem[Brion and Kumar(2005)]{BrionKumar}
Brion, M., \& Kumar, S. (2005). \emph{Frobenius splitting methods in geometry and representation theory}. Progress in Mathematics, 231. Birkh\"auser. \url{https://www-fourier.univ-grenoble-alpes.fr/~mbrion/fbook.pdf}.
\bibitem[Debarre(1994)]{Debarre}
Debarre, O. (1994). Degrees of curves in abelian varieties. \emph{Bulletin de la Soci\'et\'e Math\'ematique de France, 122}(3), 343--361. \url{https://www.math.ens.psl.eu/~debarre/25.pdf}.
\bibitem[Edixhoven et~al.(n.d.)]{EGM}
Edixhoven, B., van der Geer, G., \& Moonen, B. (n.d.). \emph{Abelian varieties: Preliminary version of the first chapters}. Author manuscript. \url{https://van-der-geer.nl/~gerard/AV.pdf}.
\bibitem[Grothendieck(n.d.)]{Grothendieck}
Grothendieck, A. (n.d.). \emph{Notes on crystalline cohomology}, cote~4, \S8.5.8, archive sheets~7--8. Institut Montpelli\'erain Alexander Grothendieck. \url{https://grothendieck.umontpellier.fr/4.pdf}.
\bibitem[Gruson and Sam(2015)]{GrusonSam}
Gruson, L., \& Sam, S. V. (2015). Alternating trilinear forms on a nine-dimensional space and degenerations of $(3,3)$-polarized abelian surfaces. \emph{Proceedings of the London Mathematical Society, 110}(3), 755--785. \url{https://doi.org/10.1112/plms/pdu050}. Author version: \url{https://arxiv.org/abs/1301.5276}.
\bibitem[Oda(1971)]{Oda}
Oda, T. (1971). Vector bundles on an elliptic curve. \emph{Nagoya Mathematical Journal, 43}, 41--72. \url{https://doi.org/10.1017/S0027763000014367}.
\bibitem[Rains and Sam(2019)]{RainsSam}
Rains, E. M., \& Sam, S. V. (2019). Vector bundles on genus~2 curves and trivectors. \emph{Algebraic Geometry, 6}(3), 328--345. \url{https://doi.org/10.14231/AG-2019-016}.
\bibitem[Sannai and Tanaka(2016)]{SannaiTanaka}
Sannai, A., \& Tanaka, H. (2016). A characterization of ordinary abelian varieties by the Frobenius push-forward of the structure sheaf. \emph{Mathematische Annalen, 366}, 1067--1087. \url{https://doi.org/10.1007/s00208-015-1352-3}.
\end{thebibliography}
\end{document}
